\ifdefined\gkver\else\def\gkver{student}\fi
\documentclass[\gkver]{gaokaozhenti}
\begin{document}

\chapter{2004年天津理综}
%% paperType: 理综物理部分

\begin{enumerate}
\item
%% number: 14
%% paperName: 2004年天津理综第 14 题 6 分
%% typeId: 1
%% type: 单选题
%% chapter: 电场
%% point: 电势能电势差电场力做功
%% method: 无
%% score: 6
%% degree: 850
%% duplicateId: 0
%% body:
在静电场中，将一电子从 A 点移到 B 点，电场力做了正功，则\xzanswer{C}

\fourchoices[answer=C]
{电场强度的方向一定是由 A 点指向 B 点}
{电场强度的方向一定是由 B 点指向 A 点}
{电子在 A 点的电势能一定比在 B 点高}
{电子在 B 点的电势能一定比在 A 点高}
%% answer: C
%% memo:
\memoanswer{
将电子从 A 点移到 B 点，电场力做正功，电势能减少，故 C 正确 D 错误；但不能确定电子是否沿电场线移动，故 AB 错误。
}

\item
%% number: 15
%% paperName: 2004年天津理综第 15 题 6 分
%% typeId: 1
%% type: 单选题
%% chapter: 气体性质
%% point: 热力学第二定律
%% method: 无
%% score: 6
%% degree: 850
%% duplicateId: 0
%% body:
下列说法正确的是\xzanswer{D}

\fourchoices[answer=D]
{热量不能由低温物体传递到高温物体}
{外界对物体做功，物体的内能必定增加}
{第二类永动机不可能制成，是因为违反了能量守恒定律}
{不可能从单一热源吸收热量并把它全部用来做功，而不引起其他变化}
%% answer: D
%% memo:
\memoanswer{
热量可以在一定条件下从低温物体传到高温物体，但不能自发地进行，故 A 错误；由热力学第一定律 $\Delta U = W + Q$，仅由 $W$ 不能确定 $\Delta U$，故 B 错误；第二类永动机不可能制成，但它没有违反能量守恒定律，而是违反了热力学第二定律，故 C 错误；D 选项是热力学第二定律的一种表述，故 D 正确。
}

\item
%% number: 16
%% paperName: 2004年天津理综第 16 题 6 分
%% typeId: 1
%% type: 单选题
%% chapter: 机械振动
%% point: 振动图象
%% method: 无
%% score: 6
%% degree: 550
%% duplicateId: 0
%% body:
公路上匀速行驶的货车受一扰动，车上货物随车厢底板上下振动但不脱离底板。一段时间内货物在竖直方向的振动可视为简谐运动，周期为 $T$。取竖直向上为正方向，以某时刻作为计时起点，即 $t = 0$，其振动图象如图所示，则\xzanswer{C}

\onepicture[w=5.0cm]{figs/16.png}
\fourchoices[answer=C]
{$t = \frac{1}{4}T$ 时，货物对车厢底板的压力最大}
{$t = \frac{1}{2}T$ 时，货物对车厢底板的压力最小}
{$t = \frac{3}{4}T$ 时，货物对车厢底板的压力最大}
{$t = \frac{3}{4}T$ 时，货物对车厢底板的压力最小}
%% answer: C
%% memo:
\memoanswer{
要使货物对车厢底板的压力最大，则车厢底板对货物的支持力最大，即要求货物向上的加速度最大，对应图象上 $t=\frac{3}{4}T$ 时刻，而 $t=\frac{1}{4}T$ 时刻压力最小，$t=\frac{1}{2}T$ 时刻处于平衡状态，故 C 正确。
}

\item
%% number: 17
%% paperName: 2004年天津理综第 17 题 6 分
%% typeId: 1
%% type: 单选题
%% chapter: 电场
%% point: 库仑定律
%% method: 无
%% score: 6
%% degree: 550
%% duplicateId: 0
%% body:
中子内有一个电荷量为 $+ \frac{2}{3}e$ 的上夸克和两个电荷量为 $- \frac{1}{3}e$ 的下夸克，一简单模型是三个夸克都在半径为 $r$ 的同一圆周上，如图\ref{2004天津理综17a}所示。下列四幅图中，能正确表示出各夸克所受静电作用力的是\xzanswer{B}

\onepicture[num=1, numstyle=arabic, labelprefix={图}, label=2004天津理综17a, w=4.6cm]{\includesvg{figs/17a}}
\fourchoices[answer=B, ispicture=true, h=2.6cm]
{figs/17b.png}
{figs/17c.png}
{figs/17d.png}
{figs/17e.png}
%% answer: B
%% memo:
\memoanswer{
对 $-\frac{1}{3}e$ 的下夸克分析，由库仑定律 $F=k\frac{Q_{1}Q_{2}}{r^{2}}$ 可知，其受到 $+\frac{2}{3}e$ 的上夸克给的引力是另一个 $-\frac{1}{3}e$ 下夸克给的斥力的 2 倍，再由平行四边形定则可知合力方向竖直向上，故 B 正确。
}

\item
%% number: 18
%% paperName: 2004年天津理综第 18 题 6 分
%% typeId: 1
%% type: 单选题
%% chapter: 交流电
%% point: 变压器
%% method: 无
%% score: 6
%% degree: 750
%% duplicateId: 0
%% body:
一台理想降压变压器从 $10\UkV$ 的线路中降压并提供 $200\UA$ 的负载电流。已知两个线圈的匝数比为 40:1，则变压器的原线圈电流、输出电压及输出功率是\xzanswer{A}

\fourchoices[answer=A]
{$5\UA$，$250\UV$，$50\UkW$}
{$5\UA$，$10\UkV$，$50\UkW$}
{$200\UA$，$250\UV$，$50\UkW$}
{$200\UA$，$10\UkV$，$2\times10^{3}\UkW$}
%% answer: A
%% memo:
\memoanswer{
由 $\frac{I_{1}}{I_{2}}=\frac{n_{2}}{n_{1}}$ 得 $I_{1}=\frac{n_{2}}{n_{1}}I_{2}=\frac{1}{40}\times200\UA=5\UA$，由 $\frac{U_{1}}{U_{2}}=\frac{n_{1}}{n_{2}}$ 得 $U_{2}=\frac{n_{2}}{n_{1}}U_{1}=\frac{1}{40}\times10\times10^{3}\UV=250\UV$，由 $P_{2}=U_{2}I_{2}=250\times200\UW=50\UkW$，故 A 正确。
}

\item
%% number: 19
%% paperName: 2004年天津理综第 19 题 6 分
%% typeId: 1
%% type: 单选题
%% chapter: 光学
%% point: 光的干涉
%% method: 无
%% score: 6
%% degree: 450
%% duplicateId: 0
%% body:
激光散斑测速是一种崭新的测速技术，它应用了光的干涉原理。用二次曝光照相所获得的“散斑对”相当于双缝干涉实验中的双缝，待测物体的速度 $v$ 与二次曝光时间间隔 $\Delta t$ 的乘积等于双缝间距。实验中可测得二次曝光时间间隔 $\Delta t$、双缝到屏之距离 $l$ 以及相邻两条亮纹间距 $\Delta x$。若所用激光波长为 $\lambda$，则该实验确定物体运动速度的表达式是\xzanswer{B}

\fourchoices[answer=B]
{$v = \frac{\lambda \Delta x}{l\Delta t}$}
{$v = \frac{l\lambda}{\Delta x\Delta t}$}
{$v = \frac{l\Delta x}{\lambda \Delta t}$}
{$v = \frac{l\Delta t}{\lambda \Delta x}$}
%% answer: B
%% memo:
\memoanswer{
双缝干涉条纹间距 $\Delta x = \frac{l}{d} \lambda$，由题意 $d = v \Delta t$，两式联立得 $v = \frac{l \lambda}{\Delta x \Delta t}$，故 B 正确。
}

\item
%% number: 20
%% paperName: 2004年天津理综第 20 题 6 分
%% typeId: 1
%% type: 单选题
%% chapter: 光学
%% point: 光子说
%% method: 无
%% score: 6
%% degree: 550
%% duplicateId: 0
%% body:
人眼对绿光最为敏感。正常人的眼睛接收到波长为 $530\Unm$ 的绿光时，只要每秒有 6 个绿光的光子射入瞳孔，眼睛就能察觉。普朗克常量为 $6.63\times10^{-34}\UJcdots$，光速为 $3.0\times10^{8}\Ums$，则人眼能察觉到绿光时所接收到的最小功率是\xzanswer{A}

\fourchoices[answer=A]
{$2.3\times10^{-18}\UW$}
{$3.8\times10^{-19}\UW$}
{$7.0\times10^{-48}\UW$}
{$1.2\times10^{-48}\UW$}
%% answer: A
%% memo:
\memoanswer{
由题意，每秒应有 6 个光子射入瞳孔，最小功率 $P = \frac{E}{t} = \frac{n h \nu}{t} = \frac{n h c}{t \lambda} = \frac{6 \times 6.6 \times 10^{-34} \times 3 \times 10^{8}}{1 \times 530 \times 10^{-9}}\UW = 2.3 \times 10^{-18}\UW$，故 A 正确。
}

\item
%% number: 21
%% paperName: 2004年天津理综第 21 题 6 分
%% typeId: 1
%% type: 单选题
%% chapter: 运动和力
%% point: 动量守恒定律
%% method: 无
%% score: 6
%% degree: 450
%% duplicateId: 0
%% body:
如图所示，光滑水平面上有大小相同的 A、B 两球在同一直线上运动。两球质量关系为 $m_{B} = 2m_{A}$，规定向右为正方向，A、B 两球的动量均为 $6\Ukgms$，运动中两球发生碰撞，碰撞后 A 球的动量增量为 $-4\Ukgms$，则\xzanswer{A}

\onepicture[w=6.5cm]{\includesvg{figs/21}}
\fourchoices[answer=A]
{左方是 A 球，碰撞后 A、B 两球速度大小之比为 2:5}
{左方是 A 球，碰撞后 A、B 两球速度大小之比为 1:10}
{右方是 A 球，碰撞后 A、B 两球速度大小之比为 2:5}
{右方是 A 球，碰撞后 A、B 两球速度大小之比为 1:10}
%% answer: A
%% memo:
\memoanswer{
由于 A、B 两球均向右运动，由题意知 $p_{A} = p_{B} = 6\Ukgms$，$m_{B} = 2m_{A}$，得 $v_{A} = 2v_{B}$，要发生碰撞，A 球应在左方；由题意知 $p_{A}' = 6 - 4 = 2\Ukgms$，$p_{B}' = 6 + 4 = 10\Ukgms$，代入 $m_{B} = 2m_{A}$，可得 $\frac{v_{A}'}{v_{B}'} = \frac{2}{5}$，故 A 正确。
}

\item
%% number: 22
%% paperName: 2004年天津理综第 22 题 15 分
%% typeId: 4
%% type: 实验题
%% chapter: 电路
%% point: 电表改装
%% method: 无
%% score: 15
%% degree: 450
%% duplicateId: 0
%% body:
现有一块 59C2 型的小量程电流表 G（表头），满偏电流为 $50\UuA$，内阻约为 $800\sim850\UO$，把它改装成 $1\UmA$、$10\UmA$ 的两量程电流表。

可供选择的器材有：
\begin{itemize}
	\item
	滑动变阻器 $R_{1}$，最大阻值 $20\UO$；
	\item
	滑动变阻器 $R_{2}$，最大阻值 $100\UkO$；
	\item
	电阻箱 $R'$，最大阻值 $9999\UO$；
	\item
	定值电阻 $R_{0}$，阻值 $1\UkO$；
	\item
	电池 $E_{1}$，电动势 $1.5\UV$；
	\item
	电池 $E_{2}$，电动势 $3.0\UV$；
	\item
	电池 $E_{3}$，电动势 $4.5\UV$（所有电池内阻均不计）；
	\item
	标准电流表 A，满偏电流 $1.5\UmA$；
	\item
	单刀单掷开关 $S_{1}$ 和 $S_{2}$，单刀双掷开关 $S_{3}$，电阻丝及导线若干。
\end{itemize}
\begin{enumerate}
	\item
	采用如图\ref{2004天津理综22a}所示电路测量表头的内阻，为提高测量精确度，选用的滑动变阻器为\_\_\_\_\_；选用的电池为\_\_\_\_\_。
	\item
	将 G 改装成两量程电流表。现有两种备选电路，示于图\ref{2004天津理综22b}和图\ref{2004天津理综22c}。图\_\_\_\_\_为合理电路，另一电路不合理的理由是\_\_\_\_\_\_\_\_\_\_\_\_\_\_\_\_\_\_\_\_\_\_\_\_\_\_\_\_\_\_。
	\item
	将改装后的电流表与标准电流表逐格进行核对（仅核对 $1\UmA$ 量程），画出所用电路图，图中待核对的电流表符号用 $A'$ 来表示。
\end{enumerate}
\onepicture[num=1, numstyle=arabic, labelprefix={图}, label=2004天津理综22a, align=right, w=5.0cm]{figs/22a.png}
\twopicture[numA=2, numB=3, numstyle=arabic, labelprefix={图}, labelA=2004天津理综22b, labelB=2004天津理综22c, align=right, hA=4.0cm, hB=4.0cm]
{figs/22b.png}
{figs/22c.png}
%% answer: 
\jdanswer{
\begin{enumerate}
	\item
	$R_{2}$（或最大阻值 $100\UkO$），$E_{3}$（或电动势 $4.5\UV$）

	\item
	2，图\ref{2004天津理综22c}电路在通电状态下，更换量程会造成两分流电阻都未并联在表头两端，以致流过表头的电流超过其满偏电流而损坏

	\item
	校对电路如图所示

\end{enumerate}
}
%% memo:
\memoanswer{
（1）根据图\ref{2004天津理综22a}的测量原理可知，只有当电路中的电流保持不变时测量的结果才较准确，由闭合电路欧姆定律 $I=\frac{E}{R+r}$ 可知，滑动变阻器的电阻 $R$ 大，电路中的电流就稳定，所以选滑动变阻器 $R_{2}$（或最大阻值 $100\UkO$）；同理，电源也应选电动势较大的，所以电源应选 $E_{3}$。

（2）选 2，当开关接到左边接线柱时，并联电阻较小，通过电流较大，所以此时对应改装电流表的较大量程，接右边接线柱时并联两个电阻，阻值较大，通过电流较小，所以此时对应量程较小。不能选 3 是因为 3 中电路在通电状态下，更换量程会造成两分流电阻都未并联在表头两端，以致流过表头的电流超过其满偏电流而损坏。

（3）为了从零到量程范围进行一一校对所以要选择滑动变阻器的分压接法，让改装表和标准表同时测量电阻 $R_{0}$ 的电流，看读数是否相同来进行校对，所以电路图如图所示。
\onepicture[h=3.8cm]{figs/22_an.png}
}

\item
%% number: 23
%% paperName: 2004年天津理综第 23 题 15 分
%% typeId: 6
%% type: 计算题
%% chapter: 磁场
%% point: 带电粒子在磁场中的运动
%% method: 无
%% score: 15
%% degree: 550
%% duplicateId: 0
%% body:
钍核 $\ce{^{230}_{90}Th}$ 发生衰变生成镭核 $\ce{^{226}_{88}Ra}$ 并放出一个粒子。设该粒子的质量为 $m$、电荷量为 $q$，它进入电势差为 $U$ 的带窄缝的平行平板电极 $S_{1}$ 和 $S_{2}$ 间电场时，其速度为 $v_{0}$，经电场加速后，沿 $Ox$ 方向进入磁感应强度为 $B$、方向垂直纸面向外的有界匀强磁场，$Ox$ 垂直平板电极 $S_{2}$，当粒子从 p 点离开磁场时，其速度方向与 $Ox$ 方位的夹角 $\theta=60^{\circ}$，如图所示，整个装置处于真空中。
\begin{enumerate}
	\item
	写出钍核衰变方程；
	\item
	求粒子在磁场中沿圆弧运动的轨道半径 $R$；
	\item
	求粒子在磁场中运动所用时间 $t$。
\end{enumerate}
\onepicture[align=right, w=5.6cm]{figs/23.png}
%% answer: 
\jdanswer{
\begin{enumerate}
	\item
	$\ce{^{230}_{90}Th -> ^{4}_{2}He + ^{226}_{88}Ra}$

	\item
	$R=\frac{m}{qB}\sqrt{\frac{2qU}{m}+v_0^2}$

	\item
	$t=\frac{\pi m}{3qB}$

\end{enumerate}
}
%% memo:
\memoanswer{
（1）钍核衰变方程为 $\ce{^{230}_{90}Th -> ^{4}_{2}He + ^{226}_{88}Ra}$ ①。

（2）设粒子离开电场时速度为 $v$，由动能定理得

$qU = \frac{1}{2}mv^2 - \frac{1}{2}mv_0^2$ ②，

粒子在磁场中，由洛伦兹力提供向心力得

$qvB = m\frac{v^2}{R}$ ③，

由②③式得

$R = \frac{m}{qB}\sqrt{\frac{2qU}{m} + v_0^2}$ ④。

（3）粒子做圆周运动的周期为

$T = \frac{2\pi R}{v} = \frac{2\pi m}{qB}$ ⑤，

粒子在磁场中运动时间为

$t = \frac{1}{6}T$ ⑥，

由⑤⑥式得

$t = \frac{\pi m}{3qB}$ ⑦。
}

\item
%% number: 24
%% paperName: 2004年天津理综第 24 题 18 分
%% typeId: 6
%% type: 计算题
%% chapter: 运动和力
%% point: 牛顿定律的应用
%% method: 无
%% score: 18
%% degree: 550
%% duplicateId: 0
%% body:
质量 $m=1.5\Ukg$ 的物块（可视为质点）在水平恒力 $F$ 作用下，从水平面上 A 点由静止开始运动，运动一段距离撤去该力，物块继续滑行 $t=2.0\Us$ 停在 B 点，已知 A、B 两点间的距离 $s=5.0\Um$，物块与水平面间的动摩擦因数 $\mu=0.20$，求恒力 $F$ 多大。
%% answer: 
\jdanswer{
$F=15\UN$
}
%% memo:
\memoanswer{
设撤去力 $F$ 前物块的位移为 $s_{1}$，撤去力 $F$ 时物块速度 $v$，物块受到的滑动摩擦力为 $F_{1}=\mu mg$，对撤去力 $F$ 后物块滑动过程应用动量定理得 $-F_{1}t=0-mv$，由运动学公式得 $s-s_{1}=\frac{v}{2}t$，对物块运动的全过程应用动能定理得 $Fs_{1}-F_{1}s=0$，由以上各式得 $F=\frac{2\mu mgs}{2s-\mu gt^{2}}$，代入数据解得 $F=15\UN$。
}

\item
%% number: 25
%% paperName: 2004年天津理综第 25 题 22 分
%% typeId: 6
%% type: 计算题
%% chapter: 磁场
%% point: 安培力
%% method: 无
%% score: 22
%% degree: 350
%% duplicateId: 0
%% body:
磁流体发电是一种新型发电方式，图\ref{2004天津理综25a}和图\ref{2004天津理综25b}是其工作原理示意图。图\ref{2004天津理综25a}中的长方体是发电导管，其中空部分的长、高、宽分别为 $l$、$a$、$b$，前后两个侧面是绝缘体，上下两个侧面是电阻可略的导体电极，这两个电极与负载电阻 $R_{1}$ 相连。整个发电导管处于图\ref{2004天津理综25b}中磁场线圈产生的匀强磁场里，磁感应强度为 $B$，方向如图所示。发电导管内有电阻率为 $\rho$ 的高温、高速电离气体沿导管向右流动，并通过专用管道导出。由于运动的电离气体受到磁场作用，产生了电动势。发电导管内电离气体流速随磁场有无而不同。设发电导管内电离气体流速处处相同，且不存在磁场时电离气体流速为 $v_{0}$，电离气体所受摩擦阻力总与流速成正比，发电导管两端的电离气体压强差 $\Delta p$ 维持恒定，求：
\begin{enumerate}
	\item
	不存在磁场时电离气体所受的摩擦阻力 $F$ 多大；
	\item
	磁流体发电机的电动势 $E$ 的大小；
	\item
	磁流体发电机发电导管的输入功率 $P$。
\end{enumerate}
\twopicture[numA=1, numB=2, numstyle=arabic, labelprefix={图}, labelA=2004天津理综25a, labelB=2004天津理综25b, align=right, hA=3.8cm, hB=3.8cm]
{figs/25a.png}
{figs/25b.png}
%% answer: 
\jdanswer{
\begin{enumerate}
	\item
	$F = ab\Delta p$

	\item
	$E = \frac{Bav_{0}}{1 + \frac{B^{2}av_{0}}{b\Delta p(R_{L} + \rho\frac{a}{bl})}}$

	\item
	$P = \frac{abv_{0}\Delta p}{1 + \frac{B^{2}av_{0}}{b\Delta p(R_{L} + \rho\frac{a}{bl})}}$

\end{enumerate}
}
%% memo:
\memoanswer{
（1）由题意知，不存在磁场时，电离气体匀速流动，由平衡条件得 $F = \Delta pS = ab\Delta p$。

（2）设磁场存在时的气体流速为 $v$，电离粒子的带电量的大小为 $q$，由于洛伦兹力的做功，正电荷聚集于上极板，负电荷聚集于下极板，使上下极板带电，产生电势差，当电势差 $E$ 稳定时，粒子受到的洛伦兹力与电场力平衡，有 $qvB = q\frac{E}{a}$，解得 $E = Bav$，故磁流体发电机的电动势为 $E = Bav$，由电阻定律知导体管的电阻为 $R = \rho \frac{a}{bl}$，电路稳定时，回路中的电流为 $I = \frac{E}{R_{L} + R} = \frac{Bav}{R_{L} + \frac{\rho a}{bl}}$，电流 $I$（电离气体）受到的安培力为 $F_{\text{安}} = B I a = \frac{B^{2} a^{2} v}{R_{L} + \frac{\rho a}{bl}}$，方向水平向左，设 $F'$ 为存在磁场时的摩擦阻力，由题意知 $\frac{F'}{F} = \frac{v}{v_{0}}$，存在磁场时，电离气体匀速流动，由平衡条件得 $ab\Delta p = F_{\text{安}} + F'$，联立解得 $E = \frac{Bav_{0}}{1 + \frac{B^{2} av_{0}}{b\Delta p(R_{L} + \frac{\rho a}{bl})}}$。

（3）磁流体发电机发电导管的输入功率为 $P = abv\Delta p$，由能量守恒得 $P = EI + F'v$，故 $P = \frac{abv_{0}\Delta p}{1 + \frac{B^{2}av_{0}}{b\Delta p(R_{L} + \frac{\rho a}{bl})}}$。
}

\end{enumerate}

\end{document}
